\section{公式与页码索引}
\label{app:equations}

本表只列正文实际引用或讨论过的公式，页码为 CMP 印刷页（351–399）。译本用页边注「原文 p.N」对应，PDF 页码为 40 页重排，不能直接换算。

\begin{center}
\small
\begin{longtable}{@{}p{2.1cm}p{6.0cm}p{1.4cm}p{4.6cm}@{}}
\toprule
编号 & 内容 & 页 & 在本书哪章 \\
\midrule
\endhead
(1.1) & 规范变换与曲率 $F_{ij}=[D_i,D_j]$ & 353 & 第 3 章 \\
(1.2) & Yang--Mills 作用量（依赖度规，被弃用） & 353 & 第 3 章 \\
(1.3) & Chern--Simons 作用量 & 354 & 第 3 章 \\
--- & 大规范变换 $L\to L+\mathrm{const}\cdot m$ & 354 & 第 3 章 \\
--- & 绕数归一化（$g^{-1}\dd g$ 的三次式） & 354 & 第 3 章 \\
(1.5) & Wilson 线 & 355 & 第 3、4 章 \\
(1.6) & 链环配分函数 & 355 & 第 3 章 \\
(2.16)(2.17) & 单圈修正 $k\to k+h^\vee$ & 359 & 第 5 章 \\
(2.25) & 大 $k$ 时 $c\to\dim G$ & 361 & 第 7 章 \\
(3.1) & $A_0=0$ 下的作用量（对时间一阶） & 367 & 第 4 章 \\
(3.2) & 泊松括号（系数 $4\pi/k$） & 367 & 第 4 章 \\
(3.3) & 高斯定律 $\epsilon^{ij}F^a_{ij}=0$ & 367 & 第 4 章 \\
(3.4) & 带标记点的高斯定律（算符值右端） & 371 & 第 4 章 \\
(3.5) & 大 $k$ 时 $\cH=\mathrm{Inv}(\otimes_i R_i)$ & 373 & 第 5 章 \\
(3.6) & $R\otimes R=\oplus_i E_i$（四点维数） & 374 & 第 5 章 \\
(4.20) & 标架反常相位 & 382 & 第 5 章 \\
(4.21) & $q=\exp\frac{2\pi\ii}{N+k}$ & 382 & 第 5、6 章 \\
(4.27)(4.28) & 手术律 & 385--386 & 第 6 章 \\
(4.36) & $N_{abc}$ 的定义（$S^2\times S^1$ 配分） & 388 & 第 6 章 \\
(4.41) & 标架反常 $Z\to Z\,e^{2\pi\ii(n-m)c/24}$ & 390 & 第 6 章 \\
(4.49) & 秩一性 & 393 & 第 6 章 \\
(4.51) & Verlinde 关键式（$S$ 对角化融合规则） & 393 & 第 6 章 \\
\bottomrule
\end{longtable}
\end{center}

三处已知译本笔误（引用前先看）：
\begin{enumerate}[leftmargin=2em]
\item 译本 §3.2 与 p.367--368 有三处把约束 (3.3) 误引为 (3.2)；读时一律按 (3.3)。
\item 原刊 (1.2) 的 Tr 归一化由脚注 1 给出；不要按 $-\frac14\int\sqrt g\,\Tr F^2$ 去补 $1/4$。
\item 早期资料曾说 p.353 漏 $1/4$，S1 证据官已判否：原刊照印，无 $1/4$、无负号。
\end{enumerate}

其它已核事实：
\begin{itemize}[leftmargin=1.6em]
\item 译本 40 页（重排），原文 49 页；公式编号沿用原文。
\item 原刊引言没有「本文解决第二个问题」这句话——那是译本导语框所加。
\item 本文是 1988 年 7 月 Swansea IAMP 报告的扩充版（p.351 脚注 *）；收稿 1988-09-12，修改 1988-09-27。
\end{itemize}
